
\section{Network Activity}

Silima is designed to operate entirely within a local area
network (LAN). No internet connection, external cloud service,
or third-party infrastructure is required for normal operation.

\subsection{Network Ports}

Silima uses the following network ports and protocols:

\begin{description}
    \item[TCP 42699:]
    HTTP and WebSocket service for student interfaces.

    \item[TCP 42614:]
    HTTP and WebSocket service for administrative interfaces.

    \item[UDP 55400:]
    RetroArch Network RetroPad controller interface.
\end{description}

\subsection{Data Privacy}

Silima does not require students to provide names, email
addresses, student identification numbers, or other
personally identifying information.

No telemetry or student metadata is intentionally collected
or stored by the application. Gameplay inputs are transmitted
to the locally hosted server for processing.

Some application resources may be cached by the student's
web browser using standard browser caching mechanisms,
as with ordinary websites. These resources are subject
to the browser's normal caching and data removal policies.

\subsection{Connection Security}

Silima currently uses HTTP and unencrypted WebSocket
connections. HTTPS and encrypted WebSocket connections
are not supported.

As a result, network traffic is not encrypted in transit.
Silima should therefore be deployed only on networks where
access to its services can be appropriately restricted.

Network administrators may use existing firewall rules
and network segmentation to control access to Silima's
services.

\subsection{Traffic Management}

Silima implements client-side message rate limiting to
reduce unnecessary network traffic.

The current implementation restricts the browser client
to approximately 45 outgoing application messages per
second, or one message every 22 milliseconds.

This restriction is implemented within the browser and
can be bypassed by a user modifying the client or using
a custom network client. It is intended to reduce
unnecessary traffic during normal operation rather than
provide protection against malicious network flooding.


\subsection{Deployment Recommendations}

Silima is intended for deployment on a trusted local area
network (LAN), such as a school or classroom network.
No internet connection is required for operation.

Silima communicates with RetroArch through a local UDP
connection on the host computer. This communication does
not need to be exposed to other devices on the network.

Student devices require access to TCP port 42699 to
interact with the application. Administrative access
is provided separately through TCP port 42614.

Network administrators are encouraged to:
\begin{itemize}
    \item Restrict administrative access to the host
          computer or other explicitly authorized devices.
    \item Allow student devices to access only the
          client-facing service.
    \item Avoid exposing Silima's services to the
          public internet.
    \item Use existing firewall rules and network
          segmentation where appropriate.
\end{itemize}

Silima currently uses unencrypted HTTP and WebSocket
connections. The application does not request personal
student information, and gameplay inputs primarily
consist of simple controller commands.

Although encryption is not currently implemented,
administrators should be aware that network traffic
may be intercepted or modified by devices with
sufficient network access.

Encrypted communication may be considered in future
releases.
